\uuid{M2ug}
\exo7id{4770}
\titre{exo7 4770}
\auteur{quercia}
\organisation{exo7}
\datecreate{2010-03-16}
\isIndication{false}
\isCorrection{false}
\chapitre{Topologie}
\sousChapitre{Topologie des espaces vectoriels normés}
\cpge{mpsi-mp2i.nombres-reels-suites.b, pcsi.nombres-reels-suites.a, ptsi.nombres-reels-suites.a, mpsi-mp2i.limites-continuite.c, pcsi.limites-continuite.c, ptsi.limites-continuite.c, bcpst1.limites-continuite.c, mp-mpi.topologie-espaces-normes.a, pc.espaces-vectoriels-normes.a, psi.espaces-vectoriels-normes.a, mp-mpi.topologie-espaces-normes.c, pc.espaces-vectoriels-normes.c, psi.espaces-vectoriels-normes.c}
\module{Analyse}
\niveau{L2}
\difficulte{}

\contenu{
\texte{
$E=C([0,1],\R)$. Soit $g\in E$. Pour tout $f\in E$ on pose
$N(f)=\sup\limits_{x\in [0,1]}\{|f(x)g(x)|\}$.
}
\begin{enumerate}
    \item \question{Donner une condition n{\'e}cessaire et suffisante sur $g$ pour que $N$ soit une
norme sur $E$.}
    \item \question{Si pour tout $x\in [0,1]$, $g(x)\ne 0$, montrer qu'alors $N$ et $\|\ \|_{\infty}$
sont des normes sur $E$ {\'e}quivalentes.}
    \item \question{D{\'e}montrer la r{\'e}ciproque de la proposition pr{\'e}c{\'e}dente.}
\end{enumerate}
}
